Para furutas evaluaciones, se recomienda registrar simultáneamente los eventos de la Arquitectura Dirigida por Eventos (EDA por sus siglas en inglés) y las métricas de rendimiento durante cada sesión de un participante, 
utilizando marcas de tiempo sincronizadas y un código anónimo común para vincular los datos. 
Al finalizar la experiencia, se podría aplicar el IPQ y asociar sus respuestas ocn los registros de esa misma sesión para saber si fue perceptible por el usuario.
Este procedimiento permitiría examinar si las variaciones de rendimiento coinciden con eventos o segmentos específicos y cómo se relacionean con las puntuaciones de presencia, sin depender de capturas realizadas en ejecuciones distintas. 